% !TEX root = executive_summary.tex
\documentclass[conference]{IEEEtran}
\IEEEoverridecommandlockouts

\usepackage{amsmath}
\usepackage{cite}
\usepackage{booktabs}
\usepackage{graphicx}
\usepackage{url}
\usepackage{hyperref}

\hypersetup{colorlinks=true,linkcolor=blue,citecolor=blue,urlcolor=blue}

\def\BibTeX{{\rm B\kern-.05em{\sc i\kern-.025em b}\kern-.08em
    T\kern-.1667em\lower.7ex\hbox{E}\kern-.125emX}}

\begin{document}

\title{Executive Summary: Validation Protocol Bias in Agricultural\\
Yield Prediction (1--2 Page Abstract)}

\author{
\IEEEauthorblockN{TAKONGMO LIONEL BLERIOT}
\IEEEauthorblockA{
\textit{Independent ML Researcher}\\
lioneltakoms25@gmail.com}
}

\maketitle

\begin{abstract}
We empirically quantify \emph{optimism bias}---performance inflation from 
spatial autocorrelation and proxy-variable leakage---when evaluating tabular 
machine learning models for crop yield prediction. Using the FAO/World Bank 
Crop Yield dataset (25{,}932 observations; 101 countries; 10 crops; 1990--2013), 
we compare five models (Global Mean, Per-Crop Mean, Ridge, Random Forest, XGBoost) 
across three validation protocols: standard Random CV, 
Leave-Country-Out (Group) CV (LCO-CV), and Expanding-Window Forward Temporal CV. 
A known data-quality issue motivates the study: the rainfall feature 
is time-invariant within every country, acting as an injective country identifier 
and injecting proxy-variable leakage. Random Forest achieves $R^2 = 0.968$ 
under Random CV but collapses to $R^2 = 0.528$ under LCO-CV (gap: $\Delta = 0.44$), 
showing no detectable difference from the naive Per-Crop Mean baseline ($R^2 = 0.531$; 
paired difference $\Delta = -0.004 \pm 0.099$, approximate 95\% CI: $[-0.050, 0.042]$). 
Permutation tests show that held-out predictions rely on rainfall 
($\Delta R^2 = 0.370 \pm 0.011$ in RF) alongside dynamic pesticides and temperature. 
Rainfall ablation indicates that proxy leakage accounts for roughly 13\% of the optimism 
gap in Random Forest (paired per-split range: $-48\%$ to $+35\%$) and 19\% in XGBoost 
($-10\%$ to $+45\%$), with the remaining gap consistent with spatial autocorrelation 
and other country feature fingerprints. Temporal CV yields intermediate nominal scores 
($R^2 \approx 0.876$), though country-identity leakage remains active across temporal 
boundaries. Structure-aware validation and simple domain baselines are essential in 
agricultural ML.
\end{abstract}

\begin{IEEEkeywords}
crop yield prediction, cross-validation bias, spatial autocorrelation, 
data leakage, random forest, XGBoost, food security
\end{IEEEkeywords}

\section{Motivation and Research Question}

Machine learning models for crop yield prediction routinely report $R^2 > 0.85\text{--}0.90$ 
under standard random splits in the literature~\cite{vanklompenburg2020}, yet policy and food 
security deployment demands generalisation to unseen geographies or future years. We ask: 
\emph{How much does random CV overestimate performance relative to structure-aware alternatives, 
and what fraction of the gap is attributable to a specific proxy leakage mechanism?}

\section{Data Audit: The Rainfall Proxy Leakage}

The \texttt{average\_rain\_fall\_mm\_per\_year} feature takes exactly one unique 
value per country across all 101 countries in the dataset~\cite{kaggle_cropyield}. 
This occurred because the raw table joined a World Bank country mean 
(\texttt{AG.LND.PRCP.MM}) without year-matching, collapsing temporal variation. 
The result is an injective 1-to-1 country identifier. Permutation importance 
tests on held-out test splits under Random CV reveal that permuting rainfall 
degrades test $R^2$ by $0.370 \pm 0.011$ in Random Forest and $0.279 \pm 0.011$ in XGBoost, 
placing it among the most influential continuous features alongside dynamic temperature 
($0.293 / 0.225$) and pesticide tonnage ($0.329 / 0.229$). Flexible tree models 
likely exploit this barcode to partition data by country.

\section{Key Results}

\begin{table}[h]
\centering
\caption{$R^2$ Summary (Full Features, mean $\pm$ SD / 20 repeated splits)}
\resizebox{\columnwidth}{!}{%
\begin{tabular}{lcccc}
\toprule
\textbf{Model} & \textbf{Random CV} & \textbf{LCO-CV} & \textbf{$\Delta$ (Gap)} \\
\midrule
Global Mean   & $-0.000\pm0.000$ & $-0.012\pm0.013$ & $-0.012$ \\
Per-Crop Mean & $0.599\pm0.006$  & $0.531\pm0.065$  & $-0.068$ \\
Ridge         & $0.631\pm0.005$  & $0.463\pm0.172$  & $-0.168$ \\
Random Forest & $0.968\pm0.003$  & $0.528\pm0.103$  & $\mathbf{-0.440}$ \\
XGBoost       & $0.945\pm0.003$  & $0.564\pm0.083$  & $\mathbf{-0.381}$ \\
\bottomrule
\end{tabular}%
}
\end{table}

\begin{table}[h]
\centering
\caption{Effect of Removing Rainfall (Paired gap change after ablation)}
\resizebox{\columnwidth}{!}{%
\begin{tabular}{lccc}
\toprule
\textbf{Model} & \textbf{Full gap} & \textbf{Ablation gap} & \textbf{Paired $\Delta$ gap} \\
\midrule
Random Forest & $0.441\pm0.104$ & $0.381\pm0.098$ & $-0.059\pm0.089$ (${\approx}13\%$) \\
XGBoost       & $0.381\pm0.084$ & $0.309\pm0.087$ & $-0.072\pm0.075$ (${\approx}19\%$) \\
\midrule
\multicolumn{4}{l}{\textit{Random CV drop: RF $-0.091\pm0.009$; XGB $-0.087\pm0.005$}} \\
\multicolumn{4}{l}{\textit{LCO-CV change: RF $-0.032\pm0.091$; XGB $-0.014\pm0.075$}} \\
\bottomrule
\end{tabular}%
}
\end{table}

\section{Conclusions}

Random CV inflates reported $R^2$ by 0.38--0.44 points for tree ensembles. 
Crucially, under LCO-CV, Random Forest ($R^2 = 0.528$) shows no detectable difference 
from predicting each crop's historical average ($R^2 = 0.531$; $p = 0.87$). 
Removing the rainfall proxy narrows the paired gap by 13--19\% on average, with the 
remaining gap reflecting spatial autocorrelation and other country feature fingerprints. 
Temporal forward evaluation yields intermediate scores ($R^2 \approx 0.876$), but 
is confounded by country identity persistence. Structure-aware validation should 
become standard in agricultural ML.

\section*{Reproducibility}
All code, models, and data pipelines are available at:\\
\url{https://github.com/lionel-tlb/crop-yield-validation-study}

\bibliographystyle{IEEEtran}
\bibliography{references}

\end{document}
